\section{进阶方向}
\label{sec:advanced}

在给出希格斯分支与库仑分支的定义之后，读者自然会问：这些定义将引向怎样的数学？有若干方向过于艰深，无法在此详述，但有兴趣的读者或许愿意进一步探索：

\subsection{稳定包络}
Aganagi\'c 与 Okounkov \cite{aganagicEllipticStable2020} 在 Maulik–Okounkov 早期工作的基础上，在每个带哈密顿 $\C^*$-作用的辛消解上，定义了被称为{\bf 椭圆稳定包络}（elliptic stable envelopes）的类。这类对象有许多例子，它们作为不同三维 $\mathcal{N}=4$ 规范理论的 $\mathfrak{M}_{\operatorname{A/B}}$ 出现。

等变稳定包络是等变椭圆上同调中的一些类，它们对应流向消解上各个 $\C^*$-不动点的\emph{刺}（thimbles；等价地说，是把实矩映射看作 Morse 函数时的稳定流形）。它们在计数几何（enumerative geometry）的研究中扮演重要角色，并被期待成为三维镜像对称的关键数学体现之一。镜像簇的椭圆稳定包络，被期待可由乘积 $\mathfrak{M}_{\operatorname{A}}\times \mathfrak{M}_{\operatorname{B}}$ 上一个自然的“Mother”类特殊化得到；在例 \ref{ex:Gr} 的情形，这一点已由 Rim\'anyi–Smirnov–Varchenko–Zhou \cite{rimanyi3dMirror2020} 证实。这一等同把物理理论中的两类参数互换：
\begin{enumerate}
    \item “质量”（masses）：它们标定 $\mathfrak{M}_{\operatorname{A}}$ 的各个消解，以及 $\mathfrak{M}_{\operatorname{B}}$ 上的 $\C^*$-作用；
    \item “Fayet–Iliopoulos（FI）参数”：它们起相反的作用，标定 $\mathfrak{M}_{\operatorname{A}}$ 上的 $\C^*$-作用与 $\mathfrak{M}_{\operatorname{B}}$ 的各个消解。
\end{enumerate}

\begin{yizhu}[稳定包络与“Mother”类]
稳定包络（stable envelope）是 Maulik–Okounkov 为计数几何引入的一类上同调类：把消解上的 $\C^*$-不动点看作实矩映射（作为 Morse 函数）的临界点，稳定包络便沿流入各不动点的刺把这些类粘合起来；椭圆情形下，它们取值于等变椭圆上同调。所谓“Mother”类，则预期是乘积 $\mathfrak{M}_{\operatorname{A}}\times \mathfrak{M}_{\operatorname{B}}$ 上的一个普适类，对两侧分别作特殊化即得各自的椭圆稳定包络——这正是镜像对称在稳定包络层面的体现。
\end{yizhu}

\subsection{范畴 \texorpdfstring{$\mathcal{O}$}{O} 的 Koszul 对偶性}

把注意力引向三维镜像对称的数学现象之一，是范畴 $\mathcal{O}$ 之间的 Koszul 对偶性。它建立在把代数 $\mathcal{A}_{\operatorname{A/B}}$ 形变量子化为非交换代数的基础上。这些形变可以理解为把 $\R^3$ 绕 $z$ 轴旋转的 $S^1$-作用纳入其中；用物理的行话说，这叫做 $\Omega$-背景。所得的代数是非交换的，因为在这一 $S^1$-作用下只有 $z$ 轴不动，而 $z$ 轴上的两个不动点无法互换位置。
\begin{enumerate}
    \item 代数 $\mathcal{A}_{\operatorname{A}}$ 的形变，通过考虑 BFN（Braverman–Finkelberg–Nakajima）空间的 $S^1$-等变同调来实现。
    \item 代数 $\mathcal{A}_{\operatorname{B}}$ 的形变，则是把 $\C[X]=\operatorname{Sym} X^*$ 换成其 Weyl 代数，后者由关系 $[x,y]=\hbar\,\Omega(x,y)$ 定义；并把“取 $\mu^{-1}(0)$ 上 $G$-不变函数”这一运算，换成哈密顿约化的非交换类比。
\end{enumerate}
范畴 $\mathcal{O}$ 是这些非交换代数上的一类特殊模构成的范畴。它可以看作稳定包络的范畴化：这里我们不再考虑流入不动点的那些刺的同调类，而是考虑在 $\mathfrak{M}_{\operatorname{A/B}}$ 的某个形变量子化上、支撑在这些刺上的模层。更多细节见 \cite[\S 3]{BLPWgco}。

Soergel 注意到，半单李代数的范畴 $\mathcal{O}$ 的主块具有一种有趣的自对偶性质：它等价于某个分次代数上不分次模构成的范畴，而在该分次代数的分次模导出范畴 $D^b(\tilde{\mathcal{O}})$ 内部，还存在原范畴的第二个“隐藏”副本。对于各种变体，例如范畴 $\mathcal{O}$ 的奇异块或抛物范畴 $\mathcal{O}$，类似的现象也会发生，但此时出现的乃是另一个范畴的副本；例如“奇异”与“抛物”这两种性质会互换。也就是说，这些范畴的分次提升是{\bf Koszul} 的，而其 Koszul 对偶是另一个（有时不同、有时相同的）同类型范畴。

正如引言中所讨论的，我们不妨设想另一个辛奇点是某个新的单李代数的幂零锥。对带 $\C^*$-作用的一般辛奇点，范畴 $\mathcal{O}$ 的定义由 Braden、Licata、Proudfoot 和第一作者在 \cite{BLPWgco} 中给出。在计算了大量例子之后，这几位作者提出了如下猜想：
\begin{conjecture}
互为镜像对偶的辛奇点（即某个三维 $\mathcal{N}=4$ 超对称规范理论的希格斯分支与库仑分支）所对应的范畴 $\mathcal{O}$，互为 Koszul 对偶。
\end{conjecture}
该猜想的一个版本（显然需要更仔细的陈述）已在 \cite{websterKoszulDuality2019} 中得到验证。这一 Koszul 对偶性的物理解释仍不确定，不过 \cite[\S 7.5]{bullimoreBoundariesMirror2016} 提出了一种解释。

\begin{yizhu}[Koszul 对偶性的原型]
Soergel 现象里的“隐藏副本”正是 Koszul 对偶的原型：在一个分次代数的分次模导出范畴之中，同时住着原范畴与它的 Koszul 对偶（对奇异块、抛物范畴等变体，出现的是另一个同类范畴的副本）。上述猜想说的正是，三维镜像对称把这一纯表示论的现象提升到辛奇点的层面：互为镜像的两个辛奇点，其范畴 $\mathcal{O}$ 应当互为 Koszul 对偶。对一般辛奇点如何定义范畴 $\mathcal{O}$，见 Braden–Licata–Proudfoot–Webster 的工作。
\end{yizhu}

\subsection{线算子}

正如局部算子（local operators）代表能在单个点上作出的观测，还有线算子（line operators）描述能沿一条直线作出的观测。研究这些算子，是把我们对三维镜像对称的研究扩展到希格斯分支与库仑分支的定义之外的一条自然途径。

我们可以用 $d$ 维扩展 TQFT（extended TQFT）的框架来描述这个范畴：它不只给一个 $(d-1)$ 维流形配一个 Hilbert 空间，而更一般地给每个余维为 $k+1$ 的流形配一个 $k$-范畴。于是，局部算子即空间 $Z(S^{d-1})$ 的描述可以推广为：把支撑在 $k$-平面上的算子等同于 $k$-范畴 $Z(S^{d-k-1})$。特别地，线算子的范畴应当由 $Z(S^{d-2})$ 给出。

对于我们关心的三维 $\mathcal{N}=4$ 理论，这些范畴可以转到 A-扭转或 B-扭转的理论之后再理解。事实上，Hilburn 与第二作者已提出了这些范畴的代数描述。这里我们只给出 $Z(S^1)$ 的一个粗略描述，其风格与 (\ref{fns:A}) 和 (\ref{fns:B}) 相仿（关于这一提案的更多讨论，见 \cite{bravermanCoulombBranches2018} 与 \cite{DGGH}）：
\begin{itemize}
	\item 在 A-扭转中，我们得到商 $N/G$ 中\emph{全纯}环路 $N((t))/G((t))$ 上的\emph{局部常值}层（即 D-模）。
	\item 在 B-扭转中，我们得到 $N/G$ 中\emph{局部常值}环路（即小环路空间）的某个版本上的\emph{全纯}层（即拟凝聚层）。
\end{itemize}
这一提案其实是推导定义 \ref{def:Coulomb} 的一条好途径：平凡线算子由从 $N[[t]]/G[[t]]$ 推前（pushforward）得到的 D-模给出，而把这个推前模的自同态朴素地算成纤维积的 Borel–Moore 同调，恰好给出定义 \ref{def:Coulomb}。

在已知的对偶对中辨识出这些范畴，是一个引人入胜却又困难的问题，人们或可称之为\emph{de Rham 三维同调镜像对称}（关于这个名称的更多背景，见下文）。Hilburn–Raskin 最近的工作 \cite{hilburnTateThesis2022} 在 $N=\C$、$G=\{1\}$ 的情形证实了这一点。

\subsection{与四维场论及 Langlands 纲领的联系}

Kapustin 与 Witten 的开创性工作 \cite{kapustinElectricmagneticDuality2007} 把几何 Langlands 对应（geometric Langlands correspondence）的一个版本，解释为具有 $\mathcal{N}=4$ 超对称的四维场论之间的物理对偶：一对 Langlands 对偶群对应的四维超对称 Yang–Mills 理论，由 S 对偶相联系。Elliott 与第二作者 \cite{elliottGeometricLanglands2018} 发展了一套数学框架，用以描述他们提案的一个变体；该变体经过（范畴化的）几何量子化，便给出几何 Langlands 对应。此外，把这一手续施于某个三维 $\mathcal{N}=4$ 理论的 A-扭转与 B-扭转，就可以得到前述的线算子范畴。

这与本文前面讨论的三维视角相衔接，因为三维 $\mathcal{N}=4$ 理论可以作为四维 $\mathcal{N}=4$ 超 Yang–Mills 理论上的边界条件出现。特别地，Gaiotto 与 Witten \cite{gaiottoSDuality2009} 研究过的这些边界条件的 $S$ 对偶，是我们寻找镜像理论最有力的工具之一。对应于线性或循环（有限或仿射 A 型）箭图的 Nakajima 箭图簇的那些理论，就是这样产生的；而这是理解第 \ref{sec:quiver-gauge-theories} 节所讨论的这些理论之对偶的最快途径。

对这一视角的数学理解仍是一个正在兴起的课题。Hilburn 与第二作者提出了整体/局部几何 Langlands 纲领与 de Rham 三维同调镜像对称陈述之间的一种新联系。在独立的工作中，Ben-Zvi、Sakellaridis 与 Venkatesh 在\term{相对 Langlands 纲领}{relative Langlands program}的框架中实现了这一物理视角，并宣布了一系列有趣的猜想，其中包括把周期与 $L$-函数的特殊值联系起来的猜想 \cite{ben-zviRelativeLanglands}。

\subsection{Betti 三维镜像对称}

代数簇的 Betti（奇异）上同调与 de Rham 上同调当然是同构的，但它们有着不同的非交换推广。二者的体现，是给定簇上局部系的模空间与平坦联络的模空间：它们通过 Riemann–Hilbert 对应解析同构，但代数上并不相同。

对我们最重要的区别在于：代数曲线的 de Rham 模空间上的复结构依赖于底曲线的复结构，而 Betti 空间则不然。类似地，Ben-Zvi 与 Nadler 提出了几何 Langlands 对应的一个“Betti”版本，它给出 Betti 模空间上拟凝聚层的自守式描述，以补充几何 Langlands 对应的“de Rham”版本 \cite{ben-zviBettiGeometric2018}。

本文前面各节的讨论同样属于 de Rham 的世界。也就是说，在我们对\term{真空模空间}{moduli spaces of vacua}的描述里，附着在 $S^2$ 上的对象依赖于这条曲线上的复结构；出现“在一个平面上全纯、在一条正交直线上常值”的结构，正是 $S^2$ 上一种退化的复结构。另一方面，在物理学家看来，与平等地对待 $\R^3$ 中所有方向相比，这一视角显得有些人为。

事实上，Kapustin 与 Witten 的提案 \cite{kapustinElectricmagneticDuality2007} 本来就是从 Betti 的视角表述的：它不依赖于曲线的复结构，而要配合通常的（de Rham）Langlands 猜想，就必须加以修改（\cite{elliottGeometricLanglands2018} 就是这样做的）。希格斯分支与库仑分支的另一个关键特征——物理上预料其存在、但从 de Rham 视角难以看到——是超凯勒度量的存在性。对于希格斯分支 $\mathfrak{M}_{\operatorname{B}}(T^*N,G)$，这类度量已用\term{超哈密顿商}{hyperhamiltonian quotient}构造出来；但在 \cite{BFN} 的框架中，很难想象如何在库仑分支上构造这样的度量。Betti 视角也许最引人入胜之处在于：由于它不依赖于复结构，它更适合扩展 TQFT 的进路。因此，正是在这个框架中，人们才能把同调镜像对称的进路推进到极致。

为使读者有个方位感，回想二维镜像对称的情形：Kontsevich 有一个惊人的洞见——镜像对称所预言的“一个理论的 A-模型与另一个理论的 B-模型之等价”，可以把握为两个三角（dg/$A_{\infty}$）范畴的等价：从辛流形的 Fukaya 范畴到复簇上凝聚层的导出范畴。这两个范畴分别是相应扭转理论的边界条件范畴，因此理论的等价可以由范畴的等价重构出来。\footnote{事实上，Kontsevich 还猜测：前述关于有理曲线的计数镜像对称，可以由这一范畴等价推出。该猜想经 Costello 进一步发展，目前是 Caldararu、Tu 及其合作者正在研究的一个活跃课题。}用扩展二维 TQFT 的话说，这就是对偶理论各自的 $Z(\operatorname{pt})$ 之间的等价。

这为三维情形提供了一个值得追寻的诱人模型。理想情况下，我们会给每个理论 $\mathcal{T}$ 的 A-扭转与 B-扭转各配一个边界条件的 2-范畴 $Z(\operatorname{pt})$，并猜想对偶理论对应的这两个范畴等价；我们不妨称之为\emph{Betti 三维同调镜像对称}。这一纲领由 Teleman \cite{telemanGaugeTheory2014} 提出，其依据是 Kapustin–Rozansky–Saulina \cite{kapustinThreedimensionalTopological2009} 针对 $\mathcal{T}(T^*N,G)$ 的提案。在交换的情形，Gammage–Hilburn–Mazel-Gee \cite{gammagePerverseSchobers2022} 与 Doan–Rezchikov \cite{DoanRezchikov2022} 最近的工作已在 A-模型 2-范畴上取得重要进展，并暗示了处理更一般情形的雄心勃勃的纲领。

% [原文 p19–p22]
